\subsubsection{IP Interrupt Enable Register (IPIER)}

Address offset: \textbf{0x08}
\newline
Reset value: 0x0000 0000

\begin{figure}[H]
    \centering
    \begin{bytefield}[bitwidth=\bflenhalfword, endianness=big]{16}
        \bitheader[lsb=16]{16-31} \\
        \bitbox{16}{Reserved}
    \end{bytefield}
\end{figure}

\begin{figure}[H]
    \centering
    \begin{bytefield}[bitwidth=\bflenhalfword, endianness=big]{16}
        \bitheader{0-15} \\
        \bitbox{7}{Reserved} &
        \bitbox{1}{SRIE} &
        \bitbox{7}{Reserved} &
        \bitbox{1}{SLIE} \\
        \bitbox[]{7}{} &
        \bitbox[]{1}{R/W} &
        \bitbox[]{7}{} &
        \bitbox[]{1}{R/W}
    \end{bytefield}
\end{figure}

\begin{itemize}
    \item \textbf{Bit 31:9 - Reserved Bits}
    \item \textbf{Bit 8 - SRIE: Sample Right Available Interrupt Enable}\\
    Wenn dieses Bit auf 1 gesetzt wird, ist der Sample Right Available Interrupt aktiviert.
    Er wird ausgelöst, wenn das SRA Bit in IPISR gesetzt wird.
    Zusätzlich müssen global die Interrupts aktiviert sein indem das GIE Bit im GIER Register gesetzt ist.
    \item \textbf{Bit 7:1 - Reserved Bits}
    \item \textbf{Bit 0 - SLIE: Sample Left Available Interrupt Enable}\\
    Wenn dieses Bit auf 1 gesetzt wird, ist der Sample Left Available Interrupt aktiviert.
    Er wird ausgelöst, wenn das SLA Bit in IPISR gesetzt wird.
    Zusätzlich müssen global die Interrupts aktiviert sein indem das GIE Bit im GIER Register gesetzt ist.
\end{itemize}